\section*{الفصل \ref{ch:b2:plant-meristems} --- النمو الخضري للنبات: المرستيمات والنمو}

\begin{solution}{exo:b2:plant-meristems:1}
\omterm{def:b2:plant-meristems:meristem}{المرستيم} جمهرة باقية من خلايا منقسمة غير متمايزة
تشتق منها أنسجة النبات. \omterm{def:b2:plant-meristems:meristem}{فالمرستيم القمي للفرع}:
الأوراق والساق والأزهار؛ \omterm{def:b2:plant-meristems:meristem}{والمرستيم القمي للجذر}: الجذر؛ \omterm{def:b2:plant-meristems:wood}{والكامبيوم
الوعائي}: \omterm{def:b2:plant-meristems:wood}{الخشب} الثانوي واللحاء الثانوي؛ \omterm{def:b2:plant-meristems:meristem}{وكامبيوم} الفلين:
\omterm{def:b2:plant-meristems:wood}{بريدرم} \omterm{def:b2:plant-meristems:wood}{القلف} (الفلين).
\end{solution}

\begin{solution}{exo:b2:plant-meristems:2}
\omterm{prop:b2:plant-meristems:ram}{قلنسوة الجذر} (الحماية، والتزليق، واستشعار الجاذبية، وخلاياها تُسلخ
عنها)؛ ثم \omterm{def:b2:plant-meristems:meristem}{المرستيم} بمركزه الساكن (الانقسام)؛ ثم \omterm{prop:b2:plant-meristems:ram}{منطقة الاستطالة}
(تطول الخلايا عشرة أضعاف إلى عشرين)؛ ثم منطقة \omterm{def:b2:cell-differentiation:differentiation}{التمايز} (الشعرات
الجذرية، \omterm{def:b2:plant-meristems:wood}{والخشب} واللحاء الناضجان). وتنشأ الجذور الجانبية من
المحيط الدائر في الجذر الناضج، مقابل أقطاب \omterm{def:b2:plant-meristems:wood}{الخشب}.
\end{solution}

\begin{solution}{exo:b2:plant-meristems:3}
قُطعت رؤوس نباتات فاصولياء: فنمت البراعم الجانبية. وفي
مجموعة ثانية تلقّى الجذمور المقطوع كتلة آغار فيها أوكسين:
فبقيت البراعم ساكنة؛ وفي مجموعة الشاهد، كتلة آغار وحدها: فنمت
البراعم. والاستنتاج: أن الأوكسين من القمة هو الذي يثبّط
البراعم الجانبية --- \omterm{prop:b2:plant-meristems:apical}{فالسيادة القمية} هرمونية.
\end{solution}

\begin{solution}{exo:b2:plant-meristems:4}
عمرها ستون سنة؛ \omterm{def:b2:plant-meristems:wood}{والخشب} العصاري الحلقات الخمس عشرة الأخيرة؛ والحي
هو \omterm{def:b2:plant-meristems:meristem}{الكامبيوم} واللحاء وخلايا الأشعة والبرانشيم في
\omterm{def:b2:plant-meristems:wood}{الخشب} العصاري \omterm{def:b2:plant-meristems:meristem}{وكامبيوم} الفلين --- أي قشرة رقيقة حول لبّ ميت.
\end{solution}

\begin{solution}{exo:b2:plant-meristems:5}
$0.1\times(P - 0.3)$: أي \qty{0.01}{h^{-1}} (بتضاعف في $\ln 2/0.01 =
\qty{69}{h}$)، و\qty{0.03}{h^{-1}} (\qty{23}{h})، و\qty{0.05}{h^{-1}}
(\qty{14}{h}). وعند \qty{0.25}{MPa} يكون الامتلاء دون عتبة
الاستسلام: فيتوقف النمو، قبل أي ذبول.
\end{solution}

\begin{solution}{exo:b2:plant-meristems:6}
% Long unbreakable runs (a chain of numerals, a Latin binomial) are one
% directional box under bidi=basic and cannot be broken; \sloppy must be in
% force when the paragraph is BROKEN, hence the \par before \endgroup.
\begingroup\sloppy
الأوراق 1--8 عند 0 و137.5 و275 و52.5 و190 و327.5 و105
و\qty{242.5}{\degree}. ومرتبةً: 0 و52.5 و105 و137.5 و190 و242.5 و275
و327.5 --- فأصغر
فجوة \qty{32.5}{\degree}. وكل ورقة تجلس في فجوة من فجوات السابقات،
فتظلل الثماني بعضها بعضًا أقل ما يمكن ويملأ
التاج قرص الضوء ملأً متساويًا.
\par\endgroup
\end{solution}

\begin{solution}{exo:b2:plant-meristems:7}
$2\pi r h\Delta r$: فعند \qty{5}{cm}، $2\pi\times 0.05\times 12\times
0.003 = \qty{0.011}{m^3}$، أي \qty{5.7}{kg}، و\qty{2.8}{kg} من الكربون؛ وعند
\qty{40}{cm}، \qty{0.090}{m^3}، أي \qty{45}{kg}، و\qty{23}{kg} من الكربون
--- أي أكثر بثماني مرات من عرض الحلقة نفسه.
\end{solution}

\begin{solution}{exo:b2:plant-meristems:8}
\emph{clv3}: لا كابح --- فمرستيمات هائلة، وأوراق وأعضاء زهرية
زائدة، وسوق مضفورة. و\emph{wus}: لا مسرّع --- \omterm{def:b2:plant-meristems:meristem}{فالمرستيم}
ينفد بعد أوراق قليلة ويُعاد تأسيسه مرارًا. والطافرة
المضاعفة: \omterm{def:b2:meiosis-heredity:vocabulary}{النمط الظاهري} \emph{wus}، فيقع WUS تحت CLV3:
إذ إن عمل CLV3 كبت WUS، ودون WUS لا شيء يُكبت.
\end{solution}

\begin{solution}{exo:b2:plant-meristems:9}
النقل: \qty{20}{h}. والانتشار: $(0.2)^{2}/(2\times 10^{-9}) =
\qty{2e7}{s}$، أي نحو 230 يومًا. فالنقل القطبي يجعل إشارة
هرمونية صالحة للاستعمال على طول نبات في غضون يوم، ويعطيها
اتجاهًا.
\end{solution}

\begin{solution}{exo:b2:plant-meristems:10}
\qty{3}{cm} في اليوم من نسيج ناضج خلاياه \qty{200}{\micro m}:
أي \num{150} خلية في الرتل في اليوم. فعلى مئة خلية منقسمة
أن تنقسم كل منها 1.5 مرة في اليوم: أي دورة نحو \qty{16}{h}.
\end{solution}

\begin{solution}{exo:b2:plant-meristems:11}
الجبريلين يطيل السلاميات؛ والنبات الذي لا يستطيع الاستجابة له
سلاميات قصيرة وساق قصيرة. والتسميد الآزوتي الغزير يجعل
قمحًا طويلًا ينمو أطول وأثقل سنبلة حتى يقع
(يرقد) ويعفن؛ أما الساق القصيرة المتيبسة فتقف، وتضع الكربون
الزائد في الحب لا في التبن، فيرتفع دليل الحصاد. وفي
البرية يُعلى على النبات القصير ويُظلَّل بجيرانه
فيخسر.
\end{solution}

\begin{solution}{exo:b2:plant-meristems:12}
تتيح المرستيمات للنبات أن يضيف أعضاء حيث ومتى تواتيه
الظروف --- أوراقًا نحو الضوء، وجذورًا نحو الماء --- وأن يتخلى
عما لا يجدي؛ وكل برعم إبطي نقطة نمو
احتياطية، فيُنجى من الرعي والنار والتقليم بالنمو من جديد؛ و
لأن المرستيمات جنينية على الدوام، ليس \omterm{def:b2:asexual-reproduction:clone}{للجينيت} عمر
محدد. أما الكلفة: فالنبات لا يستطيع أن يتحرك ولا أن يعيد تنظيم ما
بناه --- فالعضو متى وُضع بقي حيث هو، والجسم
المبني بالتراكم لا بد أن يحمل ماضيه الميت (\omterm{def:b2:plant-meristems:wood}{الخشب} القلبي) معه.
\end{solution}

\begin{solution}{pb:b2:plant-meristems:1}
\textbf{1.} $150/3 = 50$ ورقة في القمة؛ و\num{10000} في الشجرة.
\textbf{2.} 0 و137.5 و275 و52.5 و\qty{190}{\degree}.
\textbf{3.} $8\times 137.5 = 1100 = 3\times 360 + 20$؛ و$13\times 137.5
= 1787.5 = 5\times 360 - 12.5$؛ و$21\times 137.5 = 2887.5 = 8\times 360
+ 7.5$: فالورقة 22 هي أول ورقة تقع ضمن \qty{10}{\degree} من
الأولى.
\textbf{4.} كل ورقة جديدة تجلس في أوسع فجوة، فتظلل الأوراق
بعضها بعضًا أقل ما يمكن؛ والأوكسين، الذي ينزح نقله القطبي (بناقلات
PIN) جوارَ كل بدائية، هو الذي يضبط
التباعد.
\textbf{5.} $\num{10000}\times 80\,\text{سم}^{2} = \qty{80}{m^2}$: أي إن
التاج كله أوراق هذا الموسم.
\textbf{6.} كلها؛ والأوراق الأولى تُبنى من النشاء المخزون
في \omterm{def:b2:plant-meristems:wood}{الخشب} والجذور في السنة السابقة.
\textbf{7.} القمة المثلَّثة تصنع البدائيات أسرع --- اثنتين أو ثلاثًا في
\omterm{prop:b2:plant-meristems:sam}{الفترة الورقية} --- فقد تتضاعف الأوراق ومساحة الورق ثلاث مرات، لكن الفروع
تكون مضفورة، بسوق شبيهة بالشريط وترتيب أوراق
مضطرب، وكثير من الورق الزائد يظلل نفسه.
\textbf{8.} نهارًا: $0.05\times(0.55 - 0.35) = \qty{0.01}{h^{-1}}$؛ وليلًا:
$0.05\times 0.40 = \qty{0.02}{h^{-1}}$.
\textbf{9.} نهارًا: $2e^{0.12} = \qty{2.25}{cm}$، أي \qty{0.25}{cm} مضافة؛
وليلًا: $2.25e^{0.24} = \qty{2.87}{cm}$، أي \qty{0.61}{cm} مضافة.
\textbf{10.} المعدل المتوسط \qty{0.015}{h^{-1}} على مدى \qty{96}{h}: ومنه $2e^{1.44}
= \qty{8.4}{cm}$.
\textbf{11.} نهارًا: $0.05\times 0.30 = \qty{0.015}{h^{-1}}$؛ وليلًا:
$0.05\times 0.50 = \qty{0.025}{h^{-1}}$.
\textbf{12.} $0.05\times 0.05 = \qty{0.0025}{h^{-1}}$، أي ربع معدل
النهار العادي. والتعديل التناضحي يرفع محتوى الخلية من المنحلّات
حتى يعود الامتلاء، عند كمون الماء نفسه، فوق
العتبة.
\textbf{13.} نهارًا يخفض النتح كمون الماء في
الورقة والامتلاء معه، فيصغر $P - Y$؛ وليلًا يُستعاد الامتلاء.
فالنمو يحده الماء لا السكر، ساعةً بساعة.
\textbf{14.} $2\pi\times 0.10\times 8\times 0.005 = \qty{0.025}{m^3}$؛
أي \qty{12.6}{kg} جافة.
\textbf{15.} \qty{6.3}{kg} من الكربون.
\textbf{16.} $80\times 4\times 150 = \qty{48}{kg}$ من الكربون.
\textbf{17.} $6.3/48 = \qty{13}{\%}$. والمصارف الأخرى: الأوراق
نفسها، والجذور، والأفرع والأغصان، وتنفس كل
خلية حية (نحو نصف ما يُثبَّت)، والمدخرات، والأزهار
والبذور.
\textbf{18.} عند نصف قطر \qty{0.20}{m}: $2\pi\times 0.20\times 8\times
0.005 = \qty{0.050}{m^3}$، أي ضعف هذه السنة --- إذ تضاعف المحيط
الذي يعمل \omterm{def:b2:plant-meristems:meristem}{الكامبيوم} على امتداده (وستكون الشجرة أطول).
\textbf{19.} سنة جافة أو باردة؛ أو تعرية أوراق بحشرات أو صقيع
أو نار. فالسبب المناخي يضيّق الحلقات نفسها في كل شجرة في
المنطقة ويرتبط بسجلات الطقس؛ أما الأذى فيصيب شجرة واحدة
أو غابة واحدة، والصقيع أو النار يترك ندوبًا تشريحية في الحلقة.
\textbf{20.} تنمو البراعم الأقرب إلى كل قطع في غضون أيام إلى
عدة فروع؛ وفي شهر تصير الشجرة أكثف وأغزر، بفروع أكثر
لكن أقصر.
\textbf{21.} الأوكسين على القطع يحل محل إشارة القمة: فتبقى البراعم
ساكنة ولا يتبع ذلك تفرع.
\textbf{22.} كل قصّة تنزع القمم وتحرر البراعم
تحتها؛ وست مرات في السنة، يضيف كل تحرير طبقة من فروع قصيرة،
فيمتلئ السطح بها.
\textbf{23.} من براعم إبطية ساكنة على الجذمور ومن براعم عرضية
يكوّنها \omterm{def:b2:plant-meristems:meristem}{الكامبيوم}؛ وكل برعم \omterm{def:b2:plant-meristems:meristem}{مرستيم} كامل قادر
على إعادة بناء جملة الفروع، في حين أن رأس الحيوان عضو
\omterm{def:b2:vertebrate-development:induction}{منظّم} واحد لا يُعوَّض ولا نسخ احتياطية له.
\textbf{24.} وصول سيتوكينين أكثر إلى البراعم يعزز نموها
--- فريّ نبات مقلَّم يجعله يتفرع أكثر.
\textbf{25.} 50 ورقة في القمة، و\num{10000} في الشجرة؛ ومعدلا السلامية
\qty{0.01}{h^{-1}} نهارًا و\qty{0.02}{h^{-1}} ليلًا؛ \omterm{def:b2:plant-meristems:wood}{وخشب} السنة
\qty{0.025}{m^3}، أي \qty{12.6}{kg} جافة، و\qty{6.3}{kg} من الكربون.
\end{solution}
